\begin{abstract}
Se implementó un planificador de tareas MCP que construye un DAG
ponderado y distribuye trabajo con MPI por prioridad de camino restante.
La corrección del scheduler se distingue de la incertidumbre del grafo
inferido. Se comparan ejecución secuencial, FIFO y prioridad crítica.
En 54 ejecuciones controladas se conservaron los resultados; las 27
verificaciones de consultas alcanzaron exactitud, completitud y
consistencia de 1.00. La carga sensible a prioridades redujo el tiempo
\PrioritySaving{}\% frente al baseline secuencial. También se verificaron
MPI y MCP reales. La evaluación autenticada de Jev y el despliegue
multinodo siguen pendientes.
\end{abstract}

\section{Parte I Análisis algorítmico}
\subsection{Q1 Definición motivación y supuestos}
Los agentes pueden secuenciar herramientas independientes. Se busca
reducir su latencia conservando resultados, precedencias e insumos.

La entrada es un JSON con $n$ tareas, herramientas, argumentos,
referencias a resultados, dependencias explícitas, artefactos producidos
y consumidos, recursos y costos positivos. El análisis obtiene un DAG
$G=(V,E)$ con pesos $\hat w_v$ y razones auditables para cada arista.
Dados $p$ workers, se producen asignación $a(v)$, inicios $s_v$,
resultados y estados terminales. Con duraciones reales $w_v$, se minimiza
\begin{equation}
C_{\max}=\max_{v\in V}(s_v+w_v),
\end{equation}
sujeto a $s_v\geq0$, $s_v\geq s_u+w_u$ para $(u,v)\in E$, una tarea
por worker y exclusión de accesos incompatibles a recursos.
Lectura--escritura y escritura--escritura agregan precedencias;
dos lectores pueden coincidir.

El modelo es $P\mid prec\mid C_{\max}$: procesadores idénticos y tareas
indivisibles, sin expropiación ni migración. Rank 0 coordina y no cuenta
como worker. La implementación añade comunicación, serialización y
procesos aislados. Se asumen herramientas de confianza, datos
serializables y metadatos que describan insumos y efectos compartidos.
Ante fallos se exigen resultados parciales explícitos.

El problema es NP-difícil: con $E=\varnothing$, $p=2$ y pesos enteros,
decidir un makespan igual a la mitad del trabajo contiene PARTITION.
La heurística de listas permite liberación dinámica sin buscar el
óptimo combinatorio. Jev aporta decisiones estructuradas \cite{jev};
su probabilidad no demuestra que el DAG contenga todas las
dependencias reales.

\subsection{Q2 Algoritmo central y corrección}
Se combinan Kahn \cite{kahn}, camino restante y despacho al siguiente
worker libre. Se agregan aristas por \texttt{depends\_on}, referencias,
artefactos y conflictos. La línea base heurística reconoce menciones
exactas de artefactos en descripciones. Jev consulta precedencias entre
pares y costos por intervalo, sin eliminar relaciones declaradas.
Probabilidad seleccionada menor que 0.75 bloquea una dependencia;
todo ciclo se rechaza antes de ejecutar.

En orden topológico inverso se calcula
\begin{equation}
b(v)=\hat w_v+\max_{x:(v,x)\in E}b(x),\quad \max\varnothing=0.
\end{equation}
Se elige mayor $b(v)$ y se desempata por prioridad explícita y posición
en el JSON. FIFO usa orden de ingreso. Los pesos se congelan por ejecución.

\begin{figure}[!t]
\centering
\begin{minipage}{0.95\columnwidth}\small
\textbf{Algoritmo 1 Planificación dinámica}\\
1. Validar DAG; calcular $b$ e indegrees.\\
2. Encolar nodos con indegree cero.\\
3. Mientras haya tareas listas o activas:\\
\hspace*{.7em} a. Con un worker libre, extraer la mejor tarea,\\
\hspace*{1.6em} resolver argumentos y enviar trabajo.\\
\hspace*{.7em} b. Recibir la siguiente respuesta disponible.\\
\hspace*{.7em} c. Si tuvo éxito, guardar resultado y reducir\\
\hspace*{1.6em} indegrees; encolar nuevos hijos listos.\\
\hspace*{.7em} d. Si falla, reintentar solo si es idempotente\\
\hspace*{1.6em} y queda presupuesto; si no, omitir descendientes.\\
4. Detener workers; guardar reporte e historial.
\end{minipage}
\caption{Liberación por eventos, sin barreras entre niveles completos.}
\end{figure}

Invariantes: I1, el indegree residual cuenta padres sin éxito; I2,
solo se despachan tareas listas y cada worker tiene un intento activo;
I3, las referencias usan resultados exitosos; I4, un fallo definitivo
impide despachar descendientes. I1 vale inicialmente por construcción.
Cada éxito reduce una vez sus aristas; encolar al llegar a cero preserva
I1 e I2. Resolver referencias después preserva I3. Marcar descendientes
antes del siguiente despacho preserva I4. Por inducción sobre eventos
se respeta el DAG aceptado y sus exclusiones de recursos.

Todo subgrafo pendiente de un DAG finito tiene un nodo mínimo.
Si las tareas terminan o alcanzan timeout, los reintentos son finitos
y el transporte opera, se alcanzan estados terminales. La pérdida de
un rank aborta el job: no se promete continuidad.
La contribución integra dependencias auditables y ejecución MCP sobre
MPI. Listas y camino crítico son técnicas conocidas; no se propone
una nueva garantía de optimalidad.
